\section{Tiên đề ngoại diên}

\

Mặc dù cái tên có vẻ lạ, ý nghĩa của nó là một tính chất vô cùng quen thuộc của tập hợp: hai tập hợp cùng phần tử thì bằng nhau. Chúng ta phát biểu nó như sau:
$$
    \defMath{\forall x, \forall y, \left( x = y \iff  \forall z, \left( z \in x \iff z \in y \right) \right)}.
$$
Hai tập bằng nhau sẽ được thừa hưởng các tính chất cho bởi phép đồng nhất từ giải tích vị từ bậc nhất.

\begin{figure}[H]
    \begin{center}
        \begin{tikzpicture}[>=stealth]
            % Vẽ hai tập hợp
            \draw[thick] (-2,0) ellipse (1.2cm and 2cm);
            \node at (-2, 2.4) {Tập hợp $x$};

            \draw[thick] (2,0) ellipse (1.2cm and 2cm);
            \node at (2, 2.4) {Tập hợp $y$};

            % Các phần tử trong tập x
            \node[element, label=left:$z_1$] (x1) at (-2, 1) {};
            \node[element, label=left:$z_2$] (x2) at (-2, 0) {};
            \node[element, label=left:$z_3$] (x3) at (-2, -1) {};

            % Các phần tử trong tập y
            \node[element, label=right:$z_1$] (y1) at (2, 1) {};
            \node[element, label=right:$z_3$] (y2) at (2, 0) {};
            \node[element, label=right:$z_2$] (y3) at (2, -1) {};

            % Nối mũi tên hai chiều
            \draw[guide double arrow] (x1) to[bend left=20] (y1);
            \draw[guide double arrow] (x2) to (y3);
            \draw[guide double arrow] (x3) to[bend left=10] (y2);
        \end{tikzpicture}
    \end{center}
    \caption{Minh họa hai tập hợp bằng nhau $x = y = \{z_1; z_2; z_3\}$.}
\end{figure}

Trình bày bằng lời, chúng ta coi hai tập hợp là bằng nhau nếu mỗi phần tử trong một tập hợp thì đều thuộc tập hợp còn lại, và ngược lại. Trong trường hợp mà chúng ta chỉ có một chiều, thì chúng ta có định nghĩa \defText{tập con}:
$$
    \defMath{\forall x, \forall y, \left( x \subseteq y \iff  \forall z, \left( z \in x \implies z \in y \right) \right)}.
$$


\begin{figure}[H]
    \begin{center}
        \begin{tikzpicture}[>=stealth]
            % Vẽ tập hợp x
            \draw[thick] (-2,0) ellipse (1.2cm and 2cm);
            \node at (-2, 2.4) {Tập hợp $x$};

            % Vẽ tập hợp y nghiêng và to hơn
            \draw[thick, rotate around={-20:(2.5,0)}] (2.5,0) ellipse (1.6cm and 2.5cm);
            \node at (2.5, 2.8) {Tập hợp $y$};

            % Các phần tử trong tập x
            \node[element, label=left:$z_1$] (x1) at (-2, 1) {};
            \node[element, label=left:$z_2$] (x2) at (-2, 0) {};
            \node[element, label=left:$z_3$] (x3) at (-2, -1) {};

            % Các phần tử trong tập y (nhiều phần tử hơn, phân bố bên trong elip nghiêng)
            \node[element, label=right:$z_1$] (y1) at (1.8, 1) {};
            \node[element, label=right:$z_2$] (y2) at (1.8, -0.2) {};
            \node[element, label=right:$z_3$] (y3) at (2.5, -1.2) {};
            \node[element, label=right:$z_4$] (y4) at (3.2, 0.5) {};
            \node[element, label=right:$z_5$] (y5) at (2.8, 1.6) {};

            % Nối mũi tên một chiều từ x sang y
            \draw[guide arrow] (x1) to[bend left=15] (y1);
            \draw[guide arrow] (x2) to (y2);
            \draw[guide arrow] (x3) to[bend right=15] (y3);
        \end{tikzpicture}
    \end{center}
    \caption{Minh họa tập con $x \subseteq y$ hay $y \supseteq x$.}
\end{figure}
Ngoài định nghĩa cho phép $\subseteq$, chúng ta còn có các định nghĩa sau: Với mọi tập hợp $x$ và $y$,
\begin{itemize}
   \item $\defMath{x \subseteqq y \iff x \subseteq y}$,
   \item $\defMath{x \supseteq y \iff y \subseteq x}$,
   \item $\defMath{x \supseteqq y \iff x \supseteq y}$,
   \item $\defMath{x \subset y \iff x \subseteq y \land x \neq y}$\footnote{Trong nhiều tài liệu, $\subset$ có ý nghĩa tương đương với $\subseteq$.},
   \item $\defMath{x \subsetneq y \iff x \subset y}$,
   \item $\defMath{x \subsetneqq y \iff x \subsetneq y}$,
   \item $\defMath{x \supset y \iff x \supseteq y \land x \neq y}$,
   \item $\defMath{x \supsetneq y \iff x \supset y}$,
   \item $\defMath{x \supsetneqq y \iff x \supsetneq y}$.
\end{itemize}

Nếu $x$ không phải là tập con của $y$ thì sao? Chúng ta có biến đổi như sau: Với mọi tập hợp $x$ và $y$,
\begin{align*}
    x \not \subseteq y & \iff \overline{\forall z, \left( z \in x \implies z \in y \right)}; \\
    x \not \subseteq y & \iff \exists z, \overline{\left( z \in x \implies z \in y \right)}; \\
    x \not \subseteq y & \iff \exists z, \left( z \in x \land z \notin y \right).
\end{align*}
Tương tự với định nghĩa $\supseteq$, chúng ta cũng có $\forall x, \forall y, (x \not \supseteq y \iff y \not \subseteq x)$. Hình \ref{fig:li_thuyet_tap_hop:tien_de:khong_la_tap_con} minh họa cho $x \not \subseteq y$.

\begin{figure}[H]
   \begin{center}
      \begin{tikzpicture}[>=stealth]
         % Vẽ tập hợp x
         \draw[thick] (-2,0) ellipse (1.2cm and 2cm);
         \node at (-2, 2.4) {Tập hợp $x$};

         % Vẽ tập hợp y
         \draw[thick, rotate around={20:(2.5,0)}] (2.5,0) ellipse (1.6cm and 2.5cm);
         \node at (2.5, 2.8) {Tập hợp $y$};

         % Các phần tử trong tập x
         \node[element, label=left:$z_1$] (x1) at (-2, 1) {};
         \node[element, label=left:$z_2$] (x2) at (-2, 0) {};
         \node[element emphasis, label={[colorEmphasis]left:$z_3$}] (x3) at (-2, -1) {};

         % Các phần tử trong tập y
         \node[element, label=right:$z_1$] (y1) at (1.8, 1) {};
         \node[element, label=right:$z_4$] (y2) at (1.8, 0) {};
         \node[element, label=right:$z_2$] (y4) at (3.2, 0.5) {};
         \node[element, label=right:$z_5$] (y5) at (2.8, -1.0) {};

         % Nối mũi tên một chiều từ x sang y cho z_1, z_2
         \draw[guide arrow] (x1) to[bend left=15] (y1);
         \draw[guide arrow] (x2) to[bend left=15] (y4);

         % Làm nổi bật rằng z không có tương ứng trong y
         \node[colorEmphasis] (text) at (0.25, -1) {$z_3 \notin y$};
         \draw[guide arrow, colorEmphasis, densely dashed] (x3) -- (text);
      \end{tikzpicture}
   \end{center}
   \caption{Minh họa $x \not \subseteq y$ (tồn tại $z_3 \in x$ nhưng $z_3 \notin y$).}
   \label{fig:li_thuyet_tap_hop:tien_de:khong_la_tap_con}
\end{figure}

Trước khi đi đến với những tiên đề tiếp theo, chúng ta sẽ cùng tìm hiểu về các tính chất mới mà phép đồng nhất có được từ tiên đề ngoại diên. Thứ nhất, để ý rằng, do $A \iff A$ đúng với mọi $A$ từ giải tích vị từ cho nên, với
\begin{equation*}
    \forall z, z \in x \iff z \in x.
\end{equation*}
Do đó, chúng ta có \defText{tính phản xạ}:
\begin{equation*}
    \defMath{\forall x, x = x}.
\end{equation*}

Ngoài ra, phép $=$ cũng có \defText{tính đối xứng}. Để ý rằng, do tính giao hoán của phép đẳng giá, chúng ta có với mọi $x$ và $y$.
\begin{equation*}
    (\forall z, z \in x \iff z \in y) \iff (\forall z, z \in y \iff z \in x).
\end{equation*}
Sử dụng định nghĩa của phép đồng nhất để được
\begin{equation*}
    \defMath{\forall x, \forall y, x = y \iff y = x}.
\end{equation*}

Một tính chất cũng quan trọng khác, với mọi $w, x, y$, chúng ta có khẳng định sau theo định nghĩa:
\begin{equation*}
    x = w \land w = y \iff \begin{cases}
        \forall z, z \in x \iff z \in w \\
        \forall z, z \in w \iff z \in y
    \end{cases}.
\end{equation*}
Theo tính chất phân phối của lượng từ toàn cục,
\begin{align*}
    x = w \land w = y \iff\forall z, \big((z \in x \iff z \in w) \land (z \in w \iff z \in y)\big).
\end{align*}
Tính chất bắc cầu của phép đẳng giá cho
\begin{equation*}
    x = w \land w = y \implies \forall z, z \in x \iff z \in y.
\end{equation*}
Và qua đó,
\begin{equation*}
    \defMath{\forall w, \forall x, \forall y, \left( x = w \land w = y \implies x = y \right)}.
\end{equation*}
Đây là tính chất còn lại, \defText{tính bắc cầu}.

\exercise Cho $A$, $B$ và $C$ là các tập hợp. Chứng minh rằng
\begin{enumerate}
   \item $A \subseteq A$;
   \item $A = B \iff A \subseteq B \land B \subseteq A$;
   \item $A \neq B \iff \exists z, (z \in A \land z \notin B) \lor (z \notin A \land z \in B)$;
   \item $A \subseteq B \land B \subseteq C \implies A \subseteq C$.
\end{enumerate}

\solution

\setcounter{subexercise}{1}
\arabic{subexercise}. Có $\forall z, (z \in A \implies z \in A)$ hiển nhiên đúng. Do vậy, theo định nghĩa, $A \subseteq A$. Điều phải chứng minh.

\stepcounter{subexercise}
\arabic{subexercise}. Theo định nghĩa tập con,
\begin{equation*}
   \begin{cases}
      A \subseteq B \iff \forall z, (z \in A \implies z \in B) \\
      B \subseteq A \iff \forall z, (z \in B \implies z \in A)
   \end{cases}.
\end{equation*}
Do đó,
\begin{align*}
   A \subseteq B \land B \subseteq A & \iff \forall z, (z \in A \implies z \in B) \land \forall z, (z \in B \implies z \in A);    \\
   A \subseteq B \land B \subseteq A & \iff \forall z, \left( (z \in A \implies z \in B) \land (z \in B \implies z \in A)\right); \\
   A \subseteq B \land B \subseteq A & \iff \forall z, \left(z \in A \iff z \in B\right);                                         \\
   A \subseteq B \land B \subseteq A & \iff A = B.
\end{align*}
Điều phải chứng minh.

\stepcounter{subexercise}
\arabic{subexercise}. Khi biết $\overline{P \iff Q} \iff (P \land \neg Q) \lor (\neg P \land Q)$, biến đổi cơ bản cho chúng ta:
\begin{align*}
   A \neq B & \iff \overline{A = B};                                                                                \\
   A \neq B & \iff \overline{\forall z, z \in A \iff z \in B};                                                      \\
   A \neq B & \iff \exists z, \overline{z \in A \iff z \in B};                                                      \\
   A \neq B & \iff \exists z, \left( z \in A \land z \notin B \right) \lor \left( z \notin A \land z \in B \right).
\end{align*}
Điều phải chứng minh.

\stepcounter{subexercise}
\arabic{subexercise}.
\begin{equation*}
   \begin{cases}
      A \subseteq B \iff \forall z, (z \in A \implies z \in B) \\
      B \subseteq C \iff \forall z, (z \in B \implies z \in C)
   \end{cases}.
\end{equation*}
Do đó,
\begin{align*}
   A \subseteq B \land B \subseteq C & \iff \forall z, (z \in A \implies z \in B) \land \forall z, (z \in B \implies z \in C);    \\
   A \subseteq B \land B \subseteq C & \iff \forall z, \left( (z \in A \implies z \in B) \land (z \in B \implies z \in C)\right); \\
   A \subseteq B \land B \subseteq C & \implies \forall z, \left(z \in A \implies z \in C\right);                                 \\
   A \subseteq B \land B \subseteq C & \implies A \subseteq C.
\end{align*}
Điều phải chứng minh.

\exercise Cho vị từ $V()$. Chứng minh rằng 
\begin{equation*}
   \exists ! x, V(x) \implies \exists x, V(x) \land \forall y, \forall z, (V(y) \land V(z) \implies y = z).
\end{equation*}
($x$, $y$, $z$ là tập hợp.)

\solution

Giả sử $\exists ! x, V(x)$, tương đương với $\exists x, (V(x) \land \forall y, (V(y) \implies x = y))$. Cụ thể hóa $x$ thành $X$, chúng ta có $V(X)$ và $\forall y, (V(y) \implies X = y)$. Đổi hạng thức $y$ thành hạng thức $z$, được mệnh đề tương đương $\forall z, (V(z) \implies X = z)$. Bổ sung lượng từ rỗng,
\begin{equation*}
   \begin{cases}
      \forall y, \forall z, (V(y) \implies X = y) \\ 
      \forall y, \forall z, (V(z) \implies X = z)
   \end{cases}.
\end{equation*}
Do lượng từ phổ quát có tính chất phân phối đối với phép hội nên suy ra 
\begin{equation*}
   \forall y, \forall z, \begin{cases}
      V(y) \implies X = y \\
      V(z) \implies X = z
   \end{cases}.
\end{equation*}

Có $\forall y, \forall z, \begin{cases}
   V(y) \land V(z) \implies V(y) \\
   V(y) \land V(z) \implies V(z)
\end{cases}$ theo quy tắc rút gọn. Kết hợp với tính chất bắc cầu của phép kéo theo để có \begin{equation*}
   \forall y, \forall z, \begin{cases}
      V(y) \land V(z) \implies X = y \\
      V(y) \land V(z) \implies X = z
   \end{cases}.
\end{equation*}
Sử dụng tính chất phân phối một lần nữa, \begin{equation*}
   \forall y, \forall z, (V(y) \land V(z) \implies X = y \land X = z).
\end{equation*}
Do tính bắc cầu của phép đồng nhất, nên cuối cùng cũng suy ra được \begin{equation*}
   \forall y, \forall z, (V(y) \land V(z) \implies y = z).
\end{equation*}
Kết hợp với $\exists x, V(x)$ có được do $V(X)$, khi đó \begin{equation*}
   \exists x, V(x) \land \forall y, \forall z, (V(y) \land V(z) \implies y = z).
\end{equation*}
Vì vậy, chúng ta có điều phải chứng minh.


